\section{Predicting the next 30 minutes}
\subsection{What information is available when a prediction is made?}
For station $s$, issue time $t$ and threshold $q$ in millimetres, the target is
\[
Y_{s,t}=\sum_{j=1}^{6}r_{s,t+5j\,\mathrm{min}},\qquad
p_{s,t}(q)=P(Y_{s,t}>q\mid X_{s,t}).
\]
We assume that observation timestamps mark the endpoints of five-minute accumulation intervals. A training label requires all six exact future records; taking the next six available rows could incorrectly span a gap. Totals are rounded to 0.001 mm, and exceedance means strictly greater than $q$. Eligibility depends on these local windows, rather than Task 1's monthly coverage filter.

We consider two information settings for this query. A uses recent rainfall from the requested station to estimate short-term risk. B uses station and calendar information to estimate historical risk at the requested time. Our primary focus is A: recent observations provide additional information about current rainfall conditions and, as the results show, improve predictive performance. B provides a baseline without recent rainfall and can also estimate risk for future times whose recent observations are not yet available, such as tomorrow afternoon. This estimate reflects historical station and calendar patterns rather than the weather conditions of a particular future storm. The same model families are tested in both settings, allowing us to examine how much the available information matters alongside the choice of algorithm.

\textbf{A, the short-term model}, also reads recent rainfall from a supplied clean observation table. We use a pre-specified ten-minute buffer, with twelve observations ending at $t-65,t-60,\ldots,t-10$ minutes. Each record must be present and both recorded update timestamps must be no later than $t$. Missing or not-yet-updated history excludes an example instead of becoming zero rainfall.

Table~\ref{tab:updatedelay} helps explain this choice. The delay between observation and the later of the two update timestamps has a median close to ten minutes. We combine the buffer with record-level checks to account for records with longer update delays.

\begin{table}[htbp]\centering\small
\begin{tabular}{lrrrr}\toprule
Period & Observations & Median (min) & P90 (min) & Within 10 min (\%)\\\midrule
2017--2021 (training) & 28,256,245 & 9.58 & 15.08 & 68.93 \\
2017--2024 (all years) & 48,538,625 & 9.47 & 12.68 & 76.40 \\
\bottomrule\end{tabular}
\caption{Observation-to-update delay in the cleaned records. All observations have both update timestamps. P90 is the 90th percentile; the final column is the percentage with the later update no more than ten minutes after observation. These counts describe observations, not eligible prediction examples.}
\label{tab:updatedelay}
\end{table}

Figure~\ref{fig:t2timeline} shows the separation between predictors and the future label.
\plot{17_prediction_timeline}{The information used by A at issue time $t$. Markers show interval endpoints. The latest history interval ends ten minutes before issue time; the six intervals after issue time form the target.}{fig:t2timeline}
\textbf{B, the calendar model}, needs only station, time and threshold. It estimates the historical risk associated with that place and time of year or day. For example, B can answer a query about tomorrow afternoon; A can use tomorrow's recent observations only when they become available.

Both models include station ID, latitude, longitude, and sine/cosine encodings of time of day and annual phase. The annual encoding accounts for leap years; numeric year is excluded. A adds the latest three five-minute amounts, 15-, 30- and 60-minute totals, the 30-minute maximum, the hour's wet count and dry fraction, and the difference between consecutive 15-minute totals. The rainfall summaries represent recent intensity, persistence and change. Cyclic time features capture the seasonal and daily patterns without introducing a discontinuity at midnight or New Year. This gives A sixteen numeric features and B six, each alongside station ID.

Table~\ref{tab:input_features} lists the inputs in each setting. All rainfall windows end at $t-10$ minutes.
\begin{table}[htbp]\centering\small
\begin{tabular}{p{.63\linewidth}cc}\toprule
Feature group (number of numeric features) & A & B\\\midrule
Latitude and longitude (2) & Yes & Yes\\
Time of day: sine and cosine (2) & Yes & Yes\\
Annual phase: sine and cosine (2) & Yes & Yes\\
Latest three five-minute rainfall amounts (3) & Yes & No\\
Rainfall totals over 15, 30 and 60 minutes (3) & Yes & No\\
Maximum five-minute rainfall within 30 minutes (1) & Yes & No\\
Wet-interval count and dry-interval fraction over 60 minutes (2) & Yes & No\\
Difference between the latest and preceding 15-minute totals (1) & Yes & No\\\midrule
\textbf{Total numeric features} & \textbf{16} & \textbf{6}\\
Station ID (categorical; counted separately) & Yes & Yes\\\bottomrule
\end{tabular}
\caption{Input features for the two information settings. A uses recent rainfall in addition to the station and calendar inputs shared with B. For LR and GBT, the query threshold $q$ is an additional input; RF uses it when converting class probabilities to exceedance probabilities.}
\label{tab:input_features}
\end{table}


\subsection{Experimental design}
Adjacent examples share much of their history and future rainfall, so a random train/test split would mix closely related observations. We instead assign separate years to training, selection, calibration and testing, removing labels whose target end reaches the following split boundary.
\par
\begin{table}[htbp]\centering\small
\begin{tabular}{lrrp{.29\linewidth}}\toprule
Period & A examples & B examples & Purpose\\\midrule
@@SPLIT_ROWS@@
\bottomrule
\end{tabular}
\caption{Eligible station-time examples after completeness checks and boundary purging. Fitting fractions apply only to 2017--2021; validation, calibration and testing use every eligible example.}\label{tab:t2splits}
\end{table}
\par
A has fewer examples because it additionally requires complete, available history. We investigate fitting fractions of 15\%, 30\% and 100\%. Sampling uses a deterministic station--time hash with seed 2026 and does not depend on the label; the smaller samples are contained in the larger ones. The 30\% samples contain 5,608,769 examples for A and 8,233,882 for B. Even the 15\% samples contain millions of examples, but only 2,886 and 4,198 exceed 20 mm, respectively, and many windows overlap within a rainfall episode.

We compare RF, LR and GBT at each of three fitting fractions: 15\%, 30\% and 100\%, separately in A and B. Every model uses 2017--2021 fitting data and is evaluated on the same complete eligible 2024 population within its information setting. Table~\ref{tab:comparison_protocol} summarises the common protocol. Each algorithm, fraction and information setting has its own fitted model; the models are not combined into an ensemble. The RF's within-bin rainfall distributions and the historical-frequency baseline always use the full eligible training population.

\begin{table}[htbp]\centering\small
\setlength{\tabcolsep}{4pt}
\begin{tabular}{>{\raggedright\arraybackslash}p{.12\linewidth}>{\raggedright\arraybackslash}p{.18\linewidth}>{\raggedright\arraybackslash}p{.27\linewidth}>{\raggedright\arraybackslash}p{.18\linewidth}>{\raggedright\arraybackslash}p{.14\linewidth}}\toprule
Algorithm & Fitting fractions & 2022 selection & 2023 calibration & 2024 test\\\midrule
RF & 15\%, 30\%, 100\% & Reference configuration; alarm cutoff & Temperature & All eligible examples\\
LR & 15\%, 30\%, 100\% & L2 regularisation; alarm cutoff & None & All eligible examples\\
GBT & 15\%, 30\%, 100\% & Tree depth & Temperature & All eligible examples\\\bottomrule
\end{tabular}
\caption{Common model-comparison protocol in both A and B. Fitting fractions apply only to 2017--2021 anchors. RF keeps the reference 40-tree, depth-8 configuration across fractions; the saved 30\% forest is calibrated and tested using the same procedure as the other forests. Historical frequency is a full-training comparator.}
\label{tab:comparison_protocol}
\end{table}

The RF 30\% weights were already available from the earlier validation study. We reuse those standard-configuration forests without refitting and complete their 2023 calibration and 2024 evaluation. The additional one-factor RF tuning analysis remains in Appendix~\ref{sec:rfsensitivity}; it supplements the unified 2024 comparison rather than replacing a missing test cell. Model-file origins and training batches are documented in Appendix~\ref{app:experiments}.

Model selection uses equally weighted mean Brier over seven rainfall thresholds: 0, 0.2, 1, 2, 5, 10 and 20 mm. Where calibration is applied, its temperature is selected on 2023 from 0.5, 0.75, 1, 1.25, 1.5 and 2. We also evaluate 3, 7.5 and 15 mm as additional test queries between the training thresholds.

\subsection{Four approaches to the probability query}
\textbf{Distributional random forest (RF).} The forest assigns the future total to seven bins: $\{0\}$, $(0,1]$, $(1,2]$, $(2,5]$, $(5,10]$, $(10,20]$ and $(20,\infty)$~\cite{sparkml}. We convert its bin probabilities into an exceedance probability by combining them with the observed distribution of training amounts inside each bin:

\[
\widehat p(q\mid X)=\sum_{k=0}^{6}\widetilde\pi_k(X)\widehat S_k(q),
\qquad
\widetilde\pi_k=\frac{\pi_k^{1/T}}{\sum_\ell\pi_\ell^{1/T}}.
\]
Here $\pi_k(X)$ is the forest probability for bin $k$, and $\widehat S_k(q)$ is the fraction of training amounts in that bin exceeding $q$. The temperature $T$ adjusts how concentrated the bin probabilities are; $T=1$ leaves them unchanged. The resulting exceedance probability is non-increasing in the requested rainfall threshold. The within-bin distributions are shared across weather conditions and extend to the largest training total, @@MAX_TRAIN@@ mm; the model returns zero above this empirical support.

The pipeline indexes station IDs as categorical features and fits numeric median imputation; missing history is still excluded before fitting. The reference configuration has 40 trees, 25 minimum examples per node, square-root feature subsets and 0.8 subsampling. The 15\% and 100\% searches compare depths 5 and 8 and select depth 8 in both information settings; the 30\% forests use this same reference depth. @@RF_TEMPERATURE_TEXT@@

\textbf{Threshold-conditioned logistic regression (LR).} LR predicts the binary label $\mathbf1(Y>q)$ with the rainfall threshold as an additional input. Each sampled station-time example is expanded to the seven training thresholds. Station IDs are one-hot encoded and numeric inputs scaled. L2 regularisation is selected from 0.01 and 0.1 using 2022 mean Brier. Each fit allows up to 100 optimisation iterations and can stop earlier when the convergence criterion is met. LR provides a relatively simple probability model against which to assess the tree methods; its probabilities receive no separate calibration fit.

\textbf{Threshold-conditioned gradient-boosted trees (GBT).} GBT uses the same binary labels and threshold expansion as LR, but learns nonlinear relationships by fitting trees sequentially. Its pipeline indexes station IDs, imputes numeric features and assembles them with $q$. We compare depths 3 and 5 with 40 boosting iterations, step size 0.1, minimum leaf size 25, subsampling 0.8 and 128 feature bins. Every reported configuration selects depth 5. Temperature calibration is applied to the binary probability pair $(1-p,p)$: A selects $T=0.75$ and B selects $T=1$ at each fraction. Unlike the RF survival mixture, this construction does not enforce decreasing probability as $q$ increases.

\textbf{Historical frequency.} This model estimates the event frequency for each station/month/hour group. We shrink each estimate towards the global training rate using 100 pseudo-observations and use that global rate for unseen groups. The baseline captures seasonal and daily patterns without fitting a classifier, making it a useful test of whether more complex models add value, especially in setting B.

\clearpage
\section{Prediction results}
\subsection{Does recent rainfall improve the forecast?}
Our main score is the Brier score,
\[
\mathrm{BS}=\frac{1}{N}\sum_{i=1}^{N}(\widehat p_i-z_i)^2,
\qquad z_i=\mathbf1(Y_i>q),
\]
which measures probability squared error; lower is better. PR-AUC summarises the precision--recall tradeoff, while ROC-AUC measures ranking across true- and false-positive rates. We also report log loss, which penalises confident errors, and reliability, which compares predicted probabilities with observed frequencies. Spark computes AUC using 1,000-bin curve downsampling, and its PR-AUC is not average precision~\cite{sparkeval}.

We first compare the 15\% RFs on the same @@COMMON_N@@ examples, covering @@COVERAGE@@\% of B's eligible 2024 set. Each forest uses the training population defined by its information requirements. Scoring both on shared test examples gives a paired comparison of the resulting forecasts.
\par
\begin{table}[!ht]\centering\small
\begin{tabular}{rrrrrr}\toprule
$q$ (mm) & Event rate (\%) & Brier A & Brier B & PR-AUC A & PR-AUC B\\\midrule
@@COMMON_ROWS@@
\bottomrule
\end{tabular}
\caption{The 15\% RFs evaluated on identical 2024 examples. The event rate gives context for rare-event scores: a rare exceedance event can have a small Brier score even with little predictive information.}\label{tab:t2common}
\end{table}
\par
Across the seven thresholds, mean Brier falls from @@MACRO_B@@ for B to @@MACRO_A@@ for A, a @@REDUCTION@@\% relative reduction in probability squared error. The advantage is greatest for any positive rain (@@SKILL_ZERO@@\% reduction) and decreases to @@SKILL_TWENTY@@\% above 20 mm. Recent rain is therefore useful in this holdout, but the heaviest events remain difficult.
\plot{18_common_comparison}{The 15\% RF comparison on shared 2024 examples. Left: relative Brier reduction with conditional 95\% block-bootstrap intervals. Right: PR-AUC and event prevalence on a logarithmic scale. Rainfall thresholds are displayed as equally spaced categories.}{fig:t2common}
To account for overlapping labels and shared weather across stations, we form 2,000 bootstrap replicates by resampling circular blocks of seven consecutive calendar days, retaining all stations and thresholds together. The 95\% interval for the mean relative Brier reduction is @@CI_LOW@@--@@CI_HIGH@@\%, conditional on these fitted models and the observed 2024 network.

\clearpage
\subsection{How do the model families compare?}
\par
\begin{table}[!ht]\centering\small
\begin{tabular}{llrrrr}\toprule
Mode & Model & Training & Mean Brier & \shortstack{PR-AUC\\at 1 mm} & \shortstack{Log loss\\at 1 mm}\\\midrule
@@MODEL_COMPARISON_ROWS@@
\bottomrule
\end{tabular}
\caption{Model comparison on each mode's eligible 2024 population: @@COMMON_N@@ examples for A and @@B_TEST_N@@ for B. Every classifier in this table uses 100\% of its eligible training population, so the algorithm comparison holds fitting fraction fixed. Historical frequency also uses the full training population. Mean Brier averages the seven original thresholds. Table~\ref{tab:t2common} provides the comparison on shared A/B examples.}\label{tab:t2baseline}
\end{table}
\par
@@ALGORITHM_COMPARISON_TEXT@@

For B, historical frequency achieves the lowest mean Brier. GBT has slightly higher PR-AUC at 1 mm than this baseline, but its probabilities are less accurate overall. Station and calendar features describe when rain is more likely on average; they contain much less information about the development of an individual storm than recent observations do.

\subsection{How much training data is useful?}
\begin{table}[!ht]\centering\small
\begin{tabular}{llrrrr}\toprule
Mode & Model & 15\% & 30\% & 100\% & Improvement (\%)\\\midrule
@@SAMPLING_ROWS@@
\bottomrule\end{tabular}
\caption{Mean Brier on each mode's eligible 2024 population, averaged over the seven original thresholds. Improvement is the relative Brier reduction from 15\% to 100\%; positive means lower error. All nine algorithm/fraction combinations are evaluated on 2024 in each information setting. The 15\% A GBT mean comes from the retained aggregate summary; other cells use metric exports.}\label{tab:trainingsampling}
\end{table}
\plot{22_training_fraction}{Relative improvement in mean 2024 Brier over each model's 15\% result; positive values mean lower error. Points show the recorded fitting fractions, with small horizontal offsets for readability. All three algorithms have points at 15\%, 30\% and 100\%. The open A GBT marker identifies the aggregate-summary reference.}{fig:t2sampling}
@@SAMPLING_INTERPRETATION@@

The controlled RF comparison on 2022 provides a complementary check. With fitting settings held fixed, increasing the sample from 15\% to 30\% changes mean Brier by only 0.012\% in A's favour and 0.003\% against B; both conditional intervals include zero. A modest increase in tree depth is more useful for A than doubling the sample at the reference settings. Appendix~\ref{sec:rfsensitivity} gives the individual settings and paired intervals. This validation analysis explains the tuning choices; the primary sample-size comparison is the complete 2024 grid in Table~\ref{tab:trainingsampling}.

The 15\% GBT scores come from an aggregate summary for A and per-threshold metrics for B. Appendix~\ref{app:experiments} lists the retained results and model packages.

\subsection{Probability calibration and alarm decisions}
A's RF distinguishes a wider range of risks, but several high-probability bins underestimate the event frequency. Mean predicted risk is @@A_MEAN_P@@\%, against @@A_EVENT@@\% observed, and the ten-bin weighted absolute calibration error is @@A_ECE@@ percentage points. These figures use A's full eligible set, which is also the shared set. Most examples occupy low-probability bins, so the small aggregate error should be read alongside the sample counts and the deviations at higher risks.
\plot{20_reliability}{Calibration of the 15\% RFs at 1 mm on shared examples. Upper panels compare bin-average predictions with observed event rates; B is enlarged to show its two occupied bins. Lower panels show sample counts in fixed 0.1-wide probability bins across the full probability range.}{fig:t2reliability}

\plot{23_threshold_risk}{Observed event frequencies and mean predicted probabilities from RF, LR and GBT, each fitted on 30\% of training anchors, at the seven original thresholds. Each panel uses its mode's eligible 2024 examples. The vertical axis is logarithmic and thresholds are equally spaced categories. Agreement in the mean is one calibration check; the reliability plot examines variation within probability groups.}{fig:t2thresholdrisk}
@@MATCHED_TAIL_RISK_TEXT@@

Finally, a rainfall threshold and an alarm cutoff are different choices. At $q=1$ mm, changing A's probability cutoff from 0.5 to the validation-selected 0.2 raises recall from 24.7\% to 46.4\%, while precision falls from 69.3\% to 47.4\%. Cutoffs are selected by 2022 F1 after fixing the 2023 temperature, without using 2024 outcomes. An always-negative rule already achieves about 97.4\% accuracy here, which is why probability quality and the precision--recall tradeoff are more informative than accuracy.

\subsection{Queries between the training thresholds}
The additional queries at 3, 7.5 and 15 mm test behaviour inside the RF rainfall bins and between the thresholds used to fit LR and GBT. Table~\ref{tab:interior} compares RF, LR and GBT at a matched 30\% fitting fraction on each mode's complete eligible 2024 population.
\begin{table}[!ht]\centering\small
\begin{tabular}{lrrrr}\toprule
Mode & $q$ (mm) & RF Brier & LR Brier & GBT Brier\\\midrule
@@INTERIOR_ROWS@@
\bottomrule\end{tabular}
\caption{2024 results at thresholds absent from the fitting query grid. All three algorithms use 30\% of training anchors and the same eligible test population within each mode.}\label{tab:interior}
\end{table}
\FloatBarrier
RF has lower Brier than both alternatives in all six matched-fraction queries. Its global within-bin approximation therefore remains useful at these additional query points. GBT also exposes a structural issue with treating the threshold as an unconstrained predictor: a larger rainfall threshold can occasionally produce a larger exceedance probability. @@THRESHOLD_ORDER_TEXT@@ The RF mixture preserves the expected probability ordering across thresholds by construction.

Appendix~\ref{app:prediction} retains threshold-specific baseline skill, monthly skill and feature-importance diagnostics for the 15\% models. The repeated AUC values across tables can differ slightly because Spark downsamples the curves according to their partitions.
